\chapter{Εισαγωγή}
\label{ch:intro}

% ─────────────────────────────────────────────────────────────────────────────
\section{Ενεργειακή Μετάβαση και Απελευθέρωση Αγορών Ηλεκτρικής Ενέργειας}
\label{sec:energy_transition}
% ─────────────────────────────────────────────────────────────────────────────

Η ηλεκτρική ενέργεια είναι ο κινητήριος άξονας της σύγχρονης κοινωνίας: από
τη βιομηχανική παραγωγή και τις μεταφορές ως τις υπηρεσίες υγείας και τις τηλεπικοινωνίες,
ελάχιστες πτυχές της ανθρώπινης δραστηριότητας λειτουργούν σήμερα χωρίς αξιόπιστη
και οικονομικά προσιτή ηλεκτρική ενέργεια. Παράλληλα, ο ενεργειακός τομέας διέρχεται μια περίοδο αναδιάρθρωσης: η κλιματική κρίση, η ωρίμανση των Ανανεώσιμων Πηγών Ενέργειας
(Renewable Energy Sources, RES) και η ψηφιοποίηση της οικονομίας τον μετασχηματίζουν.

Αυτή η μεταμόρφωση συνοψίζεται συχνά με τον όρο \textit{τριπλή μετάβαση} ή \textit{τρία D's}:
\textit{Απορρύθμιση} (Deregulation), \textit{Αποανθρακοποίηση} (Decarbonization) και
\textit{Αποκέντρωση} (Decentralization)~\cite{weron2014electricity}. Οι τρεις αυτές δυνάμεις
λειτουργούν ταυτόχρονα και αλληλεπιδρούν, διαμορφώνοντας ένα τοπίο ολοένα
και πιο πολύπλοκο για τους συμμετέχοντες στην αγορά ηλεκτρισμού.

Η \textit{Απορρύθμιση} αναφέρεται στη μακρά διαδικασία απομάκρυνσης από τα κλασικά
μονοπωλιακά μοντέλα προς αγορές με ελεύθερο ανταγωνισμό. Το παραδοσιακό μοντέλο,
που επικρατούσε παγκοσμίως για το μεγαλύτερο μέρος του 20ού αιώνα, βασιζόταν σε κάθετα
ολοκληρωμένες επιχειρήσεις που ελέγχαν την παραγωγή, μεταφορά, διανομή και λιανική
πώληση ηλεκτρισμού. Αυτό το μοντέλο, αν και εξασφάλιζε σταθερότητα, δεν δημιουργούσε
κίνητρα για αποδοτικότητα και καινοτομία.

Η πρώτη μεγάλη ρήξη με αυτό το πρότυπο συνέβη στο Ηνωμένο Βασίλειο το 1990, όταν
τέθηκε σε λειτουργία η \textit{England and Wales Electricity Pool} και άρχισε η απελευθέρωση
της αγοράς ηλεκτρισμού~\cite{weron2014electricity}. Λίγο αργότερα, η Νορβηγία εισήγαγε το μοντέλο
της NordPool, μιας ανεξάρτητης χρηματιστηριακής αγοράς που είναι σήμερα πρότυπο
αναφοράς για όλη την Ευρώπη. Στον ευρωπαϊκό χώρο, η θεσμική αφετηρία υπήρξε η
Οδηγία 1996/92/ΕΚ για τους κοινούς κανόνες στην εσωτερική αγορά ηλεκτρισμού, η οποία
δρομολόγησε τον σταδιακό διαχωρισμό (unbundling) των δικτύων μεταφοράς από τις παραγωγικές
και εμπορικές δραστηριότητες. Η Οδηγία 2003/54/ΕΚ επέβαλε αυστηρότερο unbundling και
επιτάχυνε το άνοιγμα της αγοράς, ενώ η Οδηγία 2009/72/ΕΚ (τρίτο ενεργειακό πακέτο)
επέκτεινε τη ρύθμιση και έθεσε τις βάσεις για την ευρωπαϊκή ενοποίηση αγορών~\cite{eu_directives_electricity}.

Ένας θεμελιώδης διαχωρισμός στη μορφή των απελευθερωμένων αγορών είναι αυτός
μεταξύ \textit{Power Pool} (Αγορά Υποχρεωτικής Κοινοπραξίας) και \textit{Power Exchange}
(Χρηματιστήριο Ηλεκτρικής Ενέργειας). Στο μοντέλο Power Exchange, η αγορά λειτουργεί
ως χρηματιστήριο: οι συμμετέχοντες υποβάλλουν εθελοντικές προσφορές αγοράς και πώλησης
ενέργειας, η τιμή διαμορφώνεται από την ισορροπία προσφοράς-ζήτησης και επιτρέπονται
διεθνείς συναλλαγές και παράγωγα χρηματοοικονομικά προϊόντα. Το τελευταίο μοντέλο
επικρατεί στις σύγχρονες ευρωπαϊκές αγορές.

Η \textit{Αποανθρακοποίηση} συνιστά τον δεύτερο πυλώνα της μετάβασης και αντανακλά
την αποφασιστικότητα των ευρωπαϊκών κυβερνήσεων να μειώσουν τις εκπομπές άνθρακα
του ενεργειακού τομέα. Στόχοι όπως το «Fit for 55» της ΕΕ και το πλαίσιο της Συμφωνίας
του Παρισιού επιβάλλουν μεγάλη αύξηση του μεριδίου των ΑΠΕ και μείωση των ορυκτών
καυσίμων. Οι ΑΠΕ, και ιδίως η ηλιακή και αιολική ενέργεια, χαρακτηρίζονται από
μηδενικό οριακό κόστος παραγωγής, αλλά εισάγουν αβεβαιότητα και μεταβλητότητα στο
σύστημα λόγω της εξάρτησής τους από τις καιρικές συνθήκες. Αυτό καθιστά την πρόβλεψη
τόσο της παραγωγής όσο και της τιμής ηλεκτρισμού πιο κρίσιμη.

Η \textit{Αποκέντρωση} αφορά τη γεωγραφική και θεσμική κατανομή της παραγωγής
ενέργειας: η ανάπτυξη κατανεμημένης παραγωγής (Distributed Generation), τα νέα
εργοστάσια ΑΠΕ μικρότερης κλίμακας, τα συστήματα αποθήκευσης και οι «ενεργειακές
κοινότητες» μετατρέπουν ένα σύστημα που παλιά ήταν ιεραρχικό και κεντρικό σε ένα
δίκτυο ευέλικτων, κατανεμημένων κόμβων. Αυτή η αλλαγή θέτει νέες προκλήσεις
στη διαχείριση φορτίου και στη σταθερότητα του δικτύου.


% ─────────────────────────────────────────────────────────────────────────────
\section{Ο Ευρωπαϊκός Σχεδιασμός: Το Μοντέλο Στόχος}
\label{sec:target_model}
% ─────────────────────────────────────────────────────────────────────────────

Απέναντι στις προκλήσεις της μετάβασης, η Ευρωπαϊκή Ένωση ανέπτυξε ένα ολοκληρωμένο
ρυθμιστικό πλαίσιο για τις αγορές ηλεκτρισμού, γνωστό ως \textit{Ευρωπαϊκό Μοντέλο
Στόχος} (European Target Model). Βασισμένο στις αρχές της φιλελευθεροποίησης, της
ενοποίησης αγορών και της αποδοτικής κατανομής πόρων, το Μοντέλο Στόχος ορίζει
τέσσερα επίπεδα αγοράς που λειτουργούν σε διαδοχικές χρονικές βαθμίδες: την
\textit{Προθεσμιακή Αγορά}, την \textit{Αγορά Επόμενης Ημέρας}, την \textit{Ενδοημερήσια
Αγορά} και την \textit{Αγορά Εξισορρόπησης}~\cite{cacm2015}.

\paragraph{Προθεσμιακή Αγορά}

Η Προθεσμιακή Αγορά (Forward Market) λειτουργεί σε χρονικό ορίζοντα εβδομάδων, μηνών
ή ετών πριν από την πραγματική παράδοση της ενέργειας. Ο κεντρικός της σκοπός δεν είναι
η φυσική παράδοση ηλεκτρισμού, αλλά η παροχή εργαλείων \textit{αντιστάθμισης κινδύνου}
(hedging) στους συμμετέχοντες. Παραγωγοί και καταναλωτές μπορούν να κλειδώσουν τιμές
για μελλοντικές περιόδους, προστατεύοντας τον εαυτό τους από την τιμολογιακή αβεβαιότητα.

\paragraph{Αγορά Επόμενης Ημέρας}

Η Αγορά Επόμενης Ημέρας (\gls{dam}) είναι ο κεντρικός μηχανισμός του Μοντέλου Στόχος.
Κάθε ημέρα, συνήθως μέχρι το μεσημέρι, οι συμμετέχοντες υποβάλλουν προσφορές πώλησης
και αγοράς ηλεκτρικής ενέργειας για κάθε μια από τις 24 ώρες της επόμενης ημέρας.
Ο αλγόριθμος εκκαθάρισης αγοράς EUPHEMIA επιλύει ταυτόχρονα τις αγορές όλων των
συζευγμένων χωρών, κατανέμοντας αποδοτικά τη συνολική προσφορά και ζήτηση και λαμβάνοντας
υπόψη τις διαθέσιμες διασυνδέσεις μεταξύ ζωνών ισορροπίας~\cite{nemo_sdac}. Η τιμή
εκκαθάρισης που προκύπτει για κάθε χώρα και κάθε ώρα, η \textit{Market Clearing Price}
ή \textit{Spot Price}, είναι ο βασικός δείκτης τιμής ηλεκτρισμού και το αντικείμενο
μελέτης της παρούσας εργασίας.

Η ευρωπαϊκή ζεύξη αγορών επόμενης ημέρας υλοποιείται μέσω του μηχανισμού
\gls{sdac}, που είναι κεντρικός θεσμός του Μοντέλου Στόχος.
Μέσω του SDAC, οι εθνικές αγορές επόμενης ημέρας χωρών-μελών της ΕΕ και γειτονικών χωρών
εκκαθαρίζονται ταυτόχρονα από έναν κοινό αλγόριθμο, με αποτέλεσμα οι τιμές να συγκλίνουν
στον βαθμό που επιτρέπουν οι διασυνδέσεις~\cite{cacm2015}.

\paragraph{Ενδοημερήσια Αγορά}

Η Ενδοημερήσια Αγορά (Intraday Market) λειτουργεί μετά την εκκαθάριση της αγοράς επόμενης ημέρας και επιτρέπει στους συμμετέχοντες να αναπροσαρμόζουν τις θέσεις τους καθ' όλη
τη διάρκεια της ίδιας ημέρας, μέχρι το κλείσιμο των πυλών (gate closure). Αυτό είναι
σημαντικό για τη διαχείριση αποκλίσεων που προκύπτουν από αλλαγές στη
διαθεσιμότητα μονάδων παραγωγής, διορθωμένες προβλέψεις για παραγωγή ΑΠΕ ή απρόσμενες
διακυμάνσεις φορτίου. Η ευρωπαϊκή ζεύξη ενδοημερήσιων αγορών υλοποιείται μέσω του
\gls{sidc}~\cite{cacm2015}.

\paragraph{Αγορά Εξισορρόπησης}

Η Αγορά Εξισορρόπησης (Balancing Market) λειτουργεί σε πραγματικό χρόνο και είναι
αποκλειστικά στην αρμοδιότητα του Διαχειριστή Συστήματος Μεταφοράς (\gls{tso}). Στόχος
της είναι η διατήρηση της ισορροπίας μεταξύ παραγωγής και κατανάλωσης: αν κάποια μονάδα παραγωγής βγει εκτός λειτουργίας απρόσμενα ή το φορτίο
αποκλίνει από την πρόβλεψη, ο διαχειριστής επεμβαίνει χρησιμοποιώντας εφεδρικές
μονάδες ισορρόπησης~\cite{ipto_balancing_market}.


% ─────────────────────────────────────────────────────────────────────────────
\section{Η Ελληνική Αγορά Ηλεκτρικής Ενέργειας}
\label{sec:greek_market}
% ─────────────────────────────────────────────────────────────────────────────

Η ελληνική αγορά ηλεκτρικής ενέργειας πέρασε τις τελευταίες τρεις δεκαετίες από
δομικές μεταρρυθμίσεις που άλλαξαν τον τρόπο λειτουργίας της. Η μετάβαση από ένα
σύστημα πλήρους κρατικής ιδιοκτησίας, με τη \gls{deh} στον κεντρικό ρόλο, σε ένα
χρηματιστηριακό μοντέλο ενταγμένο στο ευρωπαϊκό πλαίσιο απαίτησε σημαντικές θεσμικές
και επιχειρησιακές αλλαγές

\paragraph{Θεσμικό Πλαίσιο: ΡΑΑΕΥ}

Η \gls{raaey} είναι η ανεξάρτητη
ρυθμιστική αρχή της ελληνικής αγοράς ηλεκτρικής ενέργειας και φυσικού αερίου~\cite{ypen_target_model_2020}.
Ο βασικός ρόλος της είναι να εξασφαλίζει τον αποτελεσματικό ανταγωνισμό στην αγορά
ενέργειας, να προστατεύει τα συμφέροντα των καταναλωτών, να εγκρίνει τιμολόγια
πρόσβασης στα δίκτυα και να εποπτεύει τη συμμόρφωση με τους κανόνες λειτουργίας.

\paragraph{Διαχείριση Συστήματος Μεταφοράς: ΑΔΜΗΕ}

Ο \gls{admhe} είναι ο
Διαχειριστής του Ελληνικού Συστήματος Μεταφοράς Ηλεκτρικής Ενέργειας. Ιδρύθηκε το
2012 στο πλαίσιο της υποχρεωτικής απόσχισης δικτύων μεταφοράς από παραγωγικές
δραστηριότητες που επέβαλε το τρίτο ευρωπαϊκό ενεργειακό πακέτο~\cite{ipto_balancing_market}.
Ο ΑΔΜΗΕ είναι υπεύθυνος για τη λειτουργία, συντήρηση και ανάπτυξη του Εθνικού
Συστήματος Μεταφοράς Ηλεκτρικής Ενέργειας (ΕΣΜΗΕ), εκτελώντας τις αρμοδιότητές του
σε θέματα ισορρόπησης του συστήματος, διαχείρισης συμφόρησης διασυνδέσεων,
παρακολούθησης της επάρκειας παραγωγής και δημοσίευσης προβλέψεων φορτίου και
ΑΠΕ για την επόμενη ημέρα. Η εκτίμηση φορτίου που δημοσιεύει ο ΑΔΜΗΕ κάθε πρωί,
πριν από την υποβολή προσφορών στην αγορά επόμενης ημέρας, είναι σημαντικό δεδομένο
εισόδου για τα μοντέλα τιμής: μια πληροφορία που είναι τυπικά διαθέσιμη πριν από
τις αποφάσεις αγοράς και δεν δημιουργεί πρόβλημα διαρροής δεδομένων.

\paragraph{Λειτουργία Αγοράς: ΔΑΠΕΕΠ και HEnEx}

Ο \gls{dapeep} ήταν ο φορέας διαχείρισης
της ελληνικής αγοράς ηλεκτρισμού για πολλά χρόνια, λειτουργώντας κυρίως ως εκκαθαριστής
του Ημερήσιου Ενεργειακού Προγραμματισμού (ΗΕΠ) στο παλαιότερο μοντέλο αγοράς Pool.
Με τη μεταρρύθμιση του 2020 και τη μετάβαση στο Μοντέλο Στόχος, οι αρμοδιότητες
αναδιαρθρώθηκαν και δημιουργήθηκε η \gls{henex}.

Το Ελληνικό Χρηματιστήριο Ενέργειας HEnEx ιδρύθηκε το 2017 και ξεκίνησε την πλήρη
εμπορική λειτουργία του τον Νοέμβριο του 2020, οπότε και η Ελλάδα εντάχθηκε στον
μηχανισμό ζεύξης αγορών επόμενης ημέρας \gls{sdac}~\cite{nemo_sdac}. Αυτή η ένταξη υπήρξε καθοριστική για την πορεία της ελληνικής αγοράς ηλεκτρισμού: από εκείνη τη στιγμή,
η ελληνική τιμή εκκαθάρισης καθορίζεται μέσω του ευρωπαϊκού αλγορίθμου EUPHEMIA,
μαζί με τις υπόλοιπες ευρωπαϊκές αγορές που συμμετέχουν στο SDAC.

\paragraph{Η Αγορά Επόμενης Ημέρας στην Ελλάδα}

Στο σύγχρονο ελληνικό μοντέλο, η διαδικασία επόμενης ημέρας ακολουθεί τυποποιημένη
διαδρομή: κάθε ημέρα πριν τις 12:00 (κεντρικής ευρωπαϊκής ώρας), όλοι οι εγγεγραμμένοι
συμμετέχοντες (παραγωγοί, προμηθευτές και έμποροι) υποβάλλουν στη HEnEx
προσφορές πώλησης και αγοράς ηλεκτρισμού για κάθε ώρα της επόμενης ημέρας. Ο αλγόριθμος
EUPHEMIA συνδυάζει τις προσφορές από όλες τις συζευγμένες αγορές, λαμβάνει υπόψη τις
μεταφορικές ικανότητες των διασυνδέσεων και βρίσκει την τιμή εκκαθάρισης (Market Clearing
Price) για κάθε ζώνη ισορροπίας και κάθε ώρα~\cite{nemo_sdac}.

Το Σχήμα~\ref{fig:market_timeline} παρουσιάζει τη χρονική ακολουθία των
αγορών ηλεκτρισμού. Οι κύριες δραστηριότητες κατανέμονται σε τρεις
ημερολογιακές ημέρες: την D-2 (δύο ημέρες πριν την παράδοση), τη D-1
(ημέρα υποβολής προσφορών) και την D (ημέρα παράδοσης). Κατά την D-1,
πριν τις 12:00 CET, υποβάλλονται οι προσφορές στη HEnEx, ο αλγόριθμος
EUPHEMIA εκτελεί την εκκαθάριση (~13:00 CET), και τα αποτελέσματα
δημοσιεύονται. Κατά την D, τρεις ενδοημερήσιοι κύκλοι (ΣΠΕΔ 1, 2, 3)
επιτρέπουν τη διόρθωση θέσεων, ενώ η Αγορά Εξισορρόπησης (ΔΕΠ) λειτουργεί
σε πραγματικό χρόνο.

\begin{figure}[htbp]
\centering
\includegraphics[width=0.85\textwidth]{xronodiagramma_agorwn}
\caption{Χρονική ακολουθία αγορών ηλεκτρισμού (D-2/D-1/D). Τα μοντέλα ML
  υποβάλλουν προβλέψεις στον ορίζοντα D-1 πριν την κλεισίδωση της DAM στις
  12:00 CET. Η Αγορά Εξισορρόπησης (ΔΕΠ) λειτουργεί σε πραγματικό χρόνο
  κατά την D.}
\label{fig:market_timeline}
\end{figure}

\paragraph{Χαρακτηριστικά της Ελληνικής Αγοράς}

Η ελληνική αγορά εμφανίζει ορισμένα ιδιαίτερα χαρακτηριστικά που την καθιστούν
ενδιαφέρουσα ως αντικείμενο μελέτης.

Πρώτον, η \textit{διείσδυση ΑΠΕ} βρίσκεται σε συνεχή ανάπτυξη. Η ηλιακή
ακτινοβολία στην Ελλάδα είναι από τις υψηλότερες στην Ευρώπη, και η εγκατεστημένη
ισχύς φωτοβολταϊκών και αιολικών πάρκων συνεχίζει να αυξάνεται. Σε πολλές ώρες κατά
τη θερινή περίοδο, οι ΑΠΕ καλύπτουν πάνω από το 70-80\% της ζήτησης, πιέζοντας τις
τιμές χαμηλά μέσω του \textit{merit order effect}: η
φθηνή παραγωγή ΑΠΕ εκτοπίζει τις ακριβότερες θερμικές μονάδες από τη στοίβα παραγωγής.

Δεύτερον, ο \textit{ρόλος του φυσικού αερίου} είναι κεντρικός. Οι μονάδες παραγωγής
με φυσικό αέριο συνιστούν συχνά τη \textit{μεθοριακή μονάδα} που καθορίζει την τελική
τιμή εκκαθάρισης, ιδιαίτερα τις ώρες αιχμής. Αυτό σημαίνει ότι οι διεθνείς τιμές
φυσικού αερίου (ιδίως το TTF ολλανδικό benchmark) ασκούν ισχυρή επίδραση στις
ελληνικές τιμές ηλεκτρισμού. Η ενεργειακή κρίση του 2021-2022, κατά την οποία οι
τιμές TTF εκτοξεύθηκαν λόγω της μειωμένης ρωσικής παροχής φυσικού αερίου, είναι
χαρακτηριστικό παράδειγμα αυτής της εξάρτησης.

Τρίτον, οι \textit{διεθνείς διασυνδέσεις} παίζουν καθοριστικό ρόλο. Η Ελλάδα
διαθέτει διασυνδέσεις με τη Βουλγαρία (μεγαλύτερη ισχύς), τη Βόρεια Μακεδονία,
την Αλβανία, την Τουρκία, μέσω της Κρήτης και νέων καλωδίων, με πρόσθετες ζώνες.
Η χρησιμοποίηση αυτών των διασυνδέσεων καθορίζεται αυτόματα από τον αλγόριθμο
EUPHEMIA, με αποτέλεσμα η ελληνική τιμή να επηρεάζεται έντονα από τις συνθήκες
στις γειτονικές αγορές.

Τέταρτον, η \textit{εποχικότητα} έχει ιδιαίτερα χαρακτηριστικά στην Ελλάδα. Λόγω
κλίματος, η θερινή περίοδος χαρακτηρίζεται από υψηλή ζήτηση ψύξης και ταυτόχρονα
υψηλή παραγωγή ηλιακής ενέργειας: δύο αντίρροπες δυνάμεις που δημιουργούν ενδιαφέροντα
πρότυπα στην τιμή. Στην περίοδο αξιολόγησης της παρούσας εργασίας, Q1 2026 (κυλιόμενο τρίμηνο
Δεκέμβριος 2025 - Φεβρουάριος 2026, και όχι το ημερολογιακό Α$'$ τρίμηνο),
ο συνδυασμός υψηλής ζήτησης θέρμανσης, μειωμένης
υδροηλεκτρικής ισχύος και αυξημένων τιμών καυσίμων οδήγησε σε έντονη μεταβλητότητα,
με μέσες τιμές που κυμάνθηκαν μεταξύ 78~€/MWh (Φεβρουάριος 2026) και 110~€/MWh
(Δεκέμβριος 2025 - Ιανουάριος 2026).


% ─────────────────────────────────────────────────────────────────────────────
\section{Πρόβλεψη Τιμών Ηλεκτρικής Ενέργειας}
\label{sec:epf_intro}
% ─────────────────────────────────────────────────────────────────────────────

\paragraph{Η Σημασία της Πρόβλεψης Τιμής}

Η πρόβλεψη τιμών ηλεκτρικής ενέργειας (Electricity Price Forecasting, \acrshort{epf}) είναι ένας από τους πιο κρίσιμους τομείς ανάλυσης σε απελευθερωμένες αγορές
ηλεκτρισμού. Το πεδίο αυτό ξεκίνησε να αναπτύσσεται στην τελευταία δεκαετία του
20ού αιώνα, παράλληλα με τις πρώτες απελευθερώσεις αγορών, και έκτοτε έχει εξελιχθεί
σε ένα ζωτικό κλάδο τόσο στην ακαδημαϊκή έρευνα όσο και στην εμπορική πρακτική.

Οι συμμετέχοντες στην αγορά αντιμετωπίζουν ημερησίως αποφάσεις που εξαρτώνται
από την εκτίμηση της μελλοντικής τιμής. Οι \textit{παραγωγοί} χρησιμοποιούν προβλέψεις
τιμής για να βελτιστοποιήσουν τις προσφορές τους στην αγορά επόμενης ημέρας: μια ακριβής
πρόβλεψη επιτρέπει τη μεγιστοποίηση του περιθωρίου κέρδους, καθορίζοντας έτσι τη στρατηγική και κατ' επέκταση την τοποθέτηση για μελλοντικές θέσεις. Οι \textit{προμηθευτές} και
\textit{έμποροι} χρησιμοποιούν τις προβλέψεις για να σχεδιάσουν αγορές ενέργειας,
να αντισταθμίσουν κίνδυνο μέσω της Προθεσμιακής Αγοράς και να διαχειριστούν χαρτοφυλάκια
παραγωγής και κατανάλωσης. Οι \textit{μεγάλοι βιομηχανικοί καταναλωτές} εκμεταλλεύονται
τις προβλέψεις για να μεταθέσουν ενεργοβόρες διεργασίες στις φθηνότερες ώρες, μειώνοντας
το ενεργειακό τους κόστος.

\paragraph{Ιδιαίτερα Χαρακτηριστικά των Τιμών Ηλεκτρισμού}

Ο Weron~\cite{weron2014electricity} στην εκτενή ανασκόπησή του κατατάσσει το EPF ανάμεσα
στα δυσκολότερα προβλήματα πρόβλεψης χρονοσειρών, εντοπίζοντας μια σειρά χαρακτηριστικών
που το διαφοροποιούν σε σχέση με άλλες χρηματοοικονομικές ή ενεργειακές χρονοσειρές:

\textit{Πολλαπλή εποχικότητα}: Οι τιμές ηλεκτρισμού εμφανίζουν ταυτόχρονα ημερήσια,
εβδομαδιαία και ετήσια μοτίβα. Σε ημερήσιο επίπεδο, παρατηρείται τυπική διπλή αιχμή
(double peak) κατά τις πρωινές και βραδινές ώρες αιχμής, ενώ τις νυχτερινές ώρες
εμφανίζονται χαμηλές τιμές. Σε εβδομαδιαίο επίπεδο, οι τιμές των Σαββατοκύριακων τείνουν
να είναι χαμηλότερες λόγω μειωμένης βιομηχανικής κατανάλωσης. Σε ετήσιο επίπεδο, η
εποχιακή δομή επηρεάζεται από τη θερμοκρασία και τη διαθεσιμότητα υδροηλεκτρικής ενέργειας.

\textit{Ακραίες τιμές (Price Spikes)}: Μια από τις πιο χαρακτηριστικές ιδιότητες των
τιμών ηλεκτρισμού είναι η εμφάνιση απότομων αιχμών, όπου η τιμή μπορεί να
πολλαπλασιαστεί σε μερικές ώρες και στη συνέχεια να επιστρέψει γρήγορα σε κανονικά
επίπεδα. Τα spikes οφείλονται συνήθως σε ταυτόχρονη εμφάνιση υψηλής ζήτησης,
χαμηλής παραγωγής ΑΠΕ και περιορισμένης διαθεσιμότητας θερμικών μονάδων~\cite{weron2014electricity}.

\textit{Μη στασιμότητα}: Αντίθετα με πολλές χρηματοοικονομικές χρονοσειρές, οι τιμές
ηλεκτρισμού δεν συμπεριφέρονται σταθερά στατιστικά διαχρονικά. Αλλαγές στη σύνθεση
της παραγωγής, στο κόστος καυσίμων, στη ζήτηση ή στις δομές της αγοράς αλλάζουν
τα στατιστικά χαρακτηριστικά της σειράς, τονίζοντας τη σημασία προσαρμοστικών προσεγγίσεων.

\textit{Μεταβλητή συσχέτιση τιμής-φορτίου}: Η τιμή συσχετίζεται γενικά θετικά με
το φορτίο, αλλά η σχέση αυτή δεν είναι σταθερή και εξαρτάται από τη σύνθεση της
παραγωγής. Σε υψηλές διεισδύσεις ΑΠΕ, το φορτίο μπορεί να αυξάνεται ενώ η τιμή
παραμένει χαμηλή ή μειώνεται, λόγω της αυξημένης παραγωγής από ΑΠΕ~\cite{weron2014electricity}.

\paragraph{Παράγοντες που Επηρεάζουν την Τιμή Ηλεκτρισμού}

Οι παράγοντες που διαμορφώνουν την τιμή ηλεκτρισμού στην αγορά επόμενης ημέρας μπορούν
να ταξινομηθούν σε τρεις βασικές κατηγορίες~\cite{weron2014electricity, lago2021}.

Από την πλευρά της \textit{προσφοράς}, ο πιο άμεσος καθοριστής της τιμής είναι η
\textit{στοίβα παραγωγής} (merit order stack): η αλληλουχία μονάδων παραγωγής
ταξινομημένων κατά αυξανόμενο οριακό κόστος. Η τελική τιμή καθορίζεται από το
οριακό κόστος της τελευταίας μονάδας που απαιτείται για να καλυφθεί η ζήτηση.

Εκτός από τη σύνθεση της παραγωγής, κρίσιμο ρόλο παίζει το \textit{κόστος καυσίμων}.
Οι τιμές φυσικού αερίου (TTF), οι τιμές λιθάνθρακα (API2) και οι τιμές δικαιωμάτων
εκπομπής CO\textsubscript{2} (EU ETS) μεταφράζονται άμεσα στα οριακά κόστη παραγωγής
και επηρεάζουν ισχυρά την τιμή εκκαθάρισης.

Από την πλευρά της \textit{ζήτησης}, το \textit{ηλεκτρικό φορτίο} είναι ο πιο άμεσος οδηγός της τιμής: υψηλή κατανάλωση σημαίνει ότι ακριβότερες μονάδες καλούνται
να καλύψουν τη ζήτηση, ανεβάζοντας τη μεθοριακή τιμή εκκαθάρισης.


% ─────────────────────────────────────────────────────────────────────────────
\section{Πρόβλεψη Ηλεκτρικού Φορτίου}
\label{sec:elf_intro}
% ─────────────────────────────────────────────────────────────────────────────

\paragraph{Η Σημασία της Πρόβλεψης Φορτίου}

Η πρόβλεψη ηλεκτρικού φορτίου (Electricity Load Forecasting, \acrshort{elf}) είναι ένας από τους παλαιότερους και πιο καλά μελετημένους τομείς στη διαχείριση ηλεκτρικών
συστημάτων. Από τις πρώτες σημαντικές εργασίες των Gross και Galiana~\cite{gross1987}
στα τέλη της δεκαετίας του '80, ο τομέας έχει συσσωρεύσει δεκαετίες εμπειρίας και
μια πλούσια βιβλιογραφία. Παρόλα αυτά, παραμένει ενεργό πεδίο έρευνας λόγω της
συνεχούς εξέλιξης των απαιτήσεων του συστήματος.

Η ακριβής πρόβλεψη φορτίου εξυπηρετεί πολλαπλές λειτουργικές ανάγκες~\cite{hong2016}.
Πρώτον, είναι απαραίτητο δεδομένο εισόδου για τον \textit{βραχυπρόθεσμο προγραμματισμό
μονάδων} (Unit Commitment). Δεύτερον, χρησιμοποιείται για τον \textit{οικονομικό
προγραμματισμό}: τη βέλτιστη κατανομή ισχύος μεταξύ μονάδων και τον σχεδιασμό αγορών
και πωλήσεων στις αγορές επόμενης ημέρας. Τρίτον, η ακριβής πρόβλεψη φορτίου είναι
απαραίτητη για τη \textit{διαχείριση αποθεμάτων ισχύος} (Spinning and Non-Spinning
Reserves).

\paragraph{Κατηγοριοποίηση κατά Ορίζοντα Πρόβλεψης}

Ένας από τους πιο ευρέως αποδεκτούς τρόπους ταξινόμησης των προβλημάτων πρόβλεψης
φορτίου είναι ο χρονικός ορίζοντας~\cite{hong2016}:

Η \textit{Πολύ Βραχυπρόθεσμη Πρόβλεψη Φορτίου} (\gls{vstlf}, έως 1 ώρα) καλύπτει
ορίζοντα από μερικά δευτερόλεπτα έως μία ώρα και χρησιμοποιείται κυρίως με αυτόματη
ρύθμιση παραγωγής.

Η \textit{Βραχυπρόθεσμη Πρόβλεψη Φορτίου} (\gls{stlf}, 1 ώρα έως 2 εβδομάδες)
συνιστά τον κρισιμότερο ορίζοντα για τις αγορές ηλεκτρισμού, καλύπτοντας τόσο
τα 24-48h για τον σχεδιασμό της αγοράς επόμενης ημέρας όσο και τον βραχυπρόθεσμο
προγραμματισμό (Unit Commitment). Η εφαρμογή που εξετάζεται στην παρούσα εργασία
εμπίπτει ακριβώς σε αυτή την κατηγορία.

Η \textit{Μεσοπρόθεσμη Πρόβλεψη Φορτίου} (\gls{mtlf}, 2 εβδομάδες έως 3 χρόνια)
χρησιμοποιείται κυρίως για τον σχεδιασμό συστημάτων συντήρησης, αγορών καυσίμων και δυναμικότητας.

Η \textit{Μακροπρόθεσμη Πρόβλεψη Φορτίου} (\gls{ltlf}, πάνω από 3 χρόνια)
καλύπτει ορίζοντα για αποφάσεις δομής: σχεδιασμός νέων σταθμών παραγωγής,
επεκτάσεις δικτύου μεταφοράς και ενεργειακός σχεδιασμός σε επίπεδο κράτους.

\paragraph{Παράγοντες που Επηρεάζουν το Φορτίο}

Σε σχέση με τους παράγοντες που επηρεάζουν τη ζήτηση ηλεκτρισμού, η βασικότερη
κατηγορία είναι τα \textit{μετεωρολογικά δεδομένα}, και ειδικότερα η θερμοκρασία
αέρα. Η σχέση μεταξύ φορτίου και θερμοκρασίας είναι μη γραμμική (U-shaped): χαμηλές
θερμοκρασίες αυξάνουν τη ζήτηση λόγω θέρμανσης, ενώ υψηλές θερμοκρασίες οδηγούν
επίσης σε υψηλό φορτίο λόγω ψύξης.

Τα \textit{ημερολογιακά χαρακτηριστικά} (ώρα ημέρας, ημέρα εβδομάδας, αργίες, εποχή,
ημέρα μήνα) εξηγούν από μόνα τους ένα μεγάλο μέρος της διακύμανσης του φορτίου.
Η εργάσιμη ημέρα έχει διαφορετικό ημερήσιο προφίλ από την Κυριακή, και μια
αργία μπορεί να μεταβάλλει τη ζήτηση συγκριτικά με μια τυπική ίδιας ημέρας
εβδομάδας.


% ─────────────────────────────────────────────────────────────────────────────
\section{Αλληλεπίδραση Τιμής και Φορτίου}
\label{sec:price_load_interaction}
% ─────────────────────────────────────────────────────────────────────────────

Τιμή ηλεκτρισμού και ηλεκτρικό φορτίο δεν είναι ανεξάρτητα μεγέθη, αλλά
συνδέονται μέσω πολύπλευρων αλληλεξαρτήσεων που λειτουργούν και στις δύο κατευθύνσεις,
αν και με διαφορετική ένταση και χρονική κλίμακα.

Στην πιο άμεση κατεύθυνση, το \textit{φορτίο επηρεάζει την τιμή μέσω του merit order}.
Σε μια αγορά που λειτουργεί με βάση την οικονομική στοίβα παραγωγής, κάθε πρόσθετη
μονάδα ζήτησης ικανοποιείται από ακριβότερες μονάδες παραγωγής, ανεβάζοντας τη
μεθοριακή τιμή εκκαθάρισης. Γι' αυτόν τον λόγο, η εκτίμηση φορτίου που δημοσιεύει
ο ΑΔΜΗΕ κάθε πρωί μπαίνει σχεδόν πάντα ως βασικό χαρακτηριστικό εισόδου σε μοντέλα
τιμής και χρησιμοποιείται ευρέως στη βιβλιογραφία του EPF.

Η αντίστροφη κατεύθυνση (\textit{επίδραση τιμής στο φορτίο}) λειτουργεί μέσω
του λεγόμενου \textit{demand response}: η δυνατότητα των καταναλωτών να ανταποκριθούν
σε σήματα τιμής μετατοπίζοντας τη ζήτησή τους χρονικά. Μεγάλοι βιομηχανικοί
καταναλωτές με ευέλικτες παραγωγικές διεργασίες (π.χ. αλουμινιοβιομηχανίες, αλυκές,
ψυγεία) μπορούν να μειώσουν τη ζήτηση κατά τις ώρες αιχμής. Ωστόσο, στον
βραχυπρόθεσμο ορίζοντα που εξετάζεται στην παρούσα εργασία, αυτή η επίδραση
παραμένει σχετικά αδύναμη.

Η \textit{κοινή μοντελοποίηση} τιμής και φορτίου ανοίγει ενδιαφέρουσες ερευνητικές
δυνατότητες. Η λεγόμενη προσέγγιση \textit{two-stage} (δύο βημάτων) εκτιμά πρώτα το
φορτίο και στη συνέχεια χρησιμοποιεί αυτή την εκτίμηση ως είσοδο στο μοντέλο
τιμής~\cite{lago2021}. Στην παρούσα εργασία, εξετάζεται εμπειρικά η επίδραση της
αντικατάστασης της επίσημης πρόβλεψης φορτίου (ΑΔΜΗΕ) από αυτόνομα εκπαιδευμένα
μοντέλα φορτίου μηχανικής μάθησης, ώστε να αξιολογηθεί αν και κατά πόσο αυτή
βελτιώνει την ακρίβεια πρόβλεψης τιμής.

Από μεθοδολογική σκοπιά, η παράλληλη αντιμετώπιση τιμής και φορτίου σε κοινό
πλαίσιο αξιολόγησης εξυπηρετεί πολλαπλούς σκοπούς. Επιτρέπει την άμεση σύγκριση
αλγόριθμων, στρατηγικών πρόβλεψης και μεθόδων βελτιστοποίησης σε δύο προβλήματα
με διαφορετικά χαρακτηριστικά: η τιμή, με την υψηλή μεταβλητότητα και τα spikes,
και το φορτίο, με τη σταθερή εποχικότητά του, αντιπροσωπεύουν δύο διαφορετικές
κατηγορίες προβλημάτων. Έτσι,τα συμπεράσματα της σύγκριση επισημαίνουν τα εγγενή
χαρακτηριστικά κάθε μεθόδου σε δύο προβλήματα με διαφορετική δυναμική.

% ─────────────────────────────────────────────────────────────────────────────
\section{Σκοπός και Ερευνητικά Ερωτήματα}
\label{sec:purpose}
\label{sec:research_questions}

Σκοπός της εργασίας είναι η αξιολόγηση και σύγκριση προσεγγίσεων βραχυπρόθεσμης πρόβλεψης για την τιμή και το ηλεκτρικό φορτίο στην ελληνική Αγορά Επόμενης Ημέρας (DAM). Η πρόβλεψη γίνεται αποκλειστικά με πληροφορία που είναι διαθέσιμη στις ρεαλιστικές συνθήκες λειτουργίας της αγοράς, σύμφωνα με το ισχύον πλαίσιο δηλαδή έως το "κλείσιμο της πύλης" (D-1, 12:00 CET).
Η σύγκριση αναπτύσσεται σε τρεις άξονες: τη στρατηγική πρόβλεψης (Teacher-Forcing, Recursive, MIMO, Direct), τον αλγόριθμο μηχανικής μάθησης και το πλαίσιο εκπαίδευσης (στατική έναντι μηνιαίας walk-forward επανεκπαίδευσης). Παράλληλα, εξετάζεται και η συμπεριφορά των μεγεθών (τιμή και φορτίο), καθώς και οι διαφορές τους ως προς τη δυσκολία πρόβλεψης.
Η αξιολόγηση γίνεται σε δύο περιόδους: πρώτα σε έναν σταθερό μήνα αναφοράς (Δεκέμβριος 2025), όπου συγκρίνεται η δομική συμπεριφορά των στρατηγικών υπό ελεγχόμενες συνθήκες. Έπειτα στο τρίμηνο Q1 2026 (Δεκέμβριος 2025-Φεβρουάριος 2026), περίοδο έντονης μεταβλητότητας και μεταβολής καθεστώτος, όπου δοκιμάζεται η ανθεκτικότητά τους, καθώς και ο βαθμός στον οποίο απαιτούν συχνή επανεκπαίδευση για να διατηρήσουν την απόδοσή τους.
Πέρα από την κατάταξη, η εργασία εξετάζει τον ρόλο των δεδομένων και της επεξεργασίας τους στην επίδοση κάθε στρατηγικής και κάθε αλγορίθμου, ερμηνεύοντας μέσω της ανάλυσης σημαντικότητας χαρακτηριστικών σε ποια πληροφορία στηρίζεται κάθε προσέγγιση.

Με βάση τον σκοπό που τέθηκε, η εργασία οργανώνεται γύρω από πέντε ερευνητικά
ερωτήματα, τα οποία διατρέχουν τη μεθοδολογία και απαντώνται συγκεντρωτικά στο
Κεφάλαιο~\ref{ch:conclusions}:

\begin{enumerate}[label=\textbf{ΕΕ\arabic*.}, leftmargin=2.2em, itemsep=4pt]
  \item Ποια στρατηγική (Recursive, MIMO, Direct), ποιος αλγόριθμος
        και ποιο πλαίσιο εκπαίδευσης (στατικό ή walk-forward) δίνουν την ακριβέστερη
         πρακτικά εφαρμόσιμη πρόβλεψη τιμής DAM και φορτίου, σε σύγκριση με το
        ιδεατό Teacher-Forcing και τα μοντέλα αναφοράς;
  \item Ποιες βασικές διαφορές εμφανίζουν τα δύο σήματα, τιμή και φορτίο, στη
        μοντελοποίηση και την πρόγνωσή τους, υπό κοινή και παράλληλη αρχιτεκτονική
        και αξιολόγηση;
  \item Πώς μεταβάλλεται η ακρίβεια των προβλέψεων ανά προβλεπόμενο μέγεθος, στρατηγική και αλγόριθμο σε περιόδους έντονης μεταβλητότητας και αλλαγής καθεστώτος;
  \item Πώς τεκμηριώνονται οι παραπάνω συγκρίσεις σε ποιοτικό, ποσοτικό και στατιστικό επίπεδο;
    
  \item Τι αντίκτυπο έχουν τα δεδομένα και ο τρόπος επεξεργασίας τους στην επίδοση
        κάθε στρατηγικής, και τι αναδεικνύει η ανάλυση σημαντικότητας χαρακτηριστικών
        ανά στρατηγική και ανά αλγόριθμο;
\end{enumerate}

Τιμή και φορτίο μοντελοποιούνται παράλληλα, με κοινή αρχιτεκτονική, ίδιο σύνολο
χαρακτηριστικών και ίδιο πλαίσιο αξιολόγησης, ώστε τα αποτελέσματά τους να είναι άμεσα
συγκρίσιμα. Οι συνεισφορές που προκύπτουν από τη διερεύνηση αυτή συνοψίζονται στο
Κεφάλαιο~\ref{ch:conclusions}.



% ─────────────────────────────────────────────────────────────────────────────
\section{Δομή της Εργασίας}
\label{sec:structure}
% ─────────────────────────────────────────────────────────────────────────────

Η εργασία οργανώνεται σε οκτώ κεφάλαια , δομημένα ώστε να
οδηγούν τον αναγνώστη βήμα προς βήμα από το θεωρητικό πλαίσιο στη μεθοδολογία, τα
αποτελέσματα και τα συμπεράσματα.

Το παρόν \textbf{Κεφάλαιο~1} παρουσίασε το ευρύτερο πλαίσιο της εργασίας: την ευρωπαϊκή
ενεργειακή μετάβαση, τη δομή και τη λειτουργία της ελληνικής αγοράς ηλεκτρισμού, τα
βασικά χαρακτηριστικά των προβλημάτων πρόβλεψης τιμής και φορτίου, καθώς και τον
σκοπό και τους άξονες διερεύνησης της εργασίας.

Το \textbf{Κεφάλαιο~\ref{ch:literature}} τοποθετεί την εργασία στο ευρύτερο ερευνητικό
πλαίσιο μέσα από βιβλιογραφική ανασκόπηση. Εξετάζεται η εξέλιξη των αλγόριθμων
Machine Learning (ML) για EPF και ELF, οι στρατηγικές πολυβήματης πρόβλεψης, η χρήση
τεχνικών Scheduled Sampling και η πρακτική της walk-forward αξιολόγησης στη σχετική
βιβλιογραφία.

Τα \textbf{Κεφάλαια~\ref{ch:data} και \ref{ch:methods}} συνιστούν τον κορμό της
μεθοδολογίας. Το Κεφάλαιο~\ref{ch:data} περιγράφει τις πηγές δεδομένων, τη διαδικασία
συλλογής, τον καθαρισμό, τη μηχανική χαρακτηριστικών (feature engineering) και τα μέτρα
αποφυγής διαρροής δεδομένων. Το Κεφάλαιο~\ref{ch:methods} παρουσιάζει αναλυτικά τους
πέντε αλγόριθμους μηχανικής μάθησης, τις τέσσερις στρατηγικές πρόβλεψης, τη
βελτιστοποίηση υπερπαραμέτρων με Optuna και το πλαίσιο walk-forward μηνιαίας
επανεκπαίδευσης.

Το \textbf{Κεφάλαιο~\ref{ch:dashboard}} παρουσιάζει το διαδραστικό εργαλείο
οπτικοποίησης (dashboard) που αναπτύχθηκε για την εξερεύνηση και σύγκριση όλων των
μοντέλων, στρατηγικών και χρονικών περιόδων αξιολόγησης.

Τα \textbf{Κεφάλαια~\ref{ch:experiments} και~\ref{ch:fi}}
συνιστούν τον πυρήνα της πειραματικής ανάλυσης. Το Κεφάλαιο~\ref{ch:experiments}
παρουσιάζει τα αριθμητικά αποτελέσματα για τιμή και φορτίο, τόσο για τη στατική
αξιολόγηση σε Δεκέμβριο 2025 όσο και για την walk-forward περίοδο Q1 2026, ανά
στρατηγική και αλγόριθμο, μαζί με ποιοτική ανάλυση των σφαλμάτων (χρονική δομή,
κατανομή, μεροληψία) και τη στατιστική επικύρωση των κατατάξεων μέσω του ελέγχου
Diebold-Mariano. Το Κεφάλαιο~\ref{ch:fi}
αναλύει τη σημαντικότητα των χαρακτηριστικών ανά στρατηγική, εξηγώντας σε ποιες
μεταβλητές στηρίζεται κάθε μοντέλο.

Η εργασία ολοκληρώνεται με το \textbf{Κεφάλαιο~\ref{ch:conclusions}}, το οποίο
αποτιμά κατά πόσο επιτεύχθηκε ο σκοπός της εργασίας, συνοψίζει τα κύρια ευρήματα και
τις συνεισφορές, οριοθετεί τους περιορισμούς της προσέγγισης και προτείνει
κατευθύνσεις για μελλοντική έρευνα.

